\cvsection{Expériences professionnelles}

\begin{cventries}
  \cventry
    {Stagiaire ingénieur de recherche en IA}
    {INRIA}
    {Bordeaux, France}
    {Mars 2026 -- Septembre 2026}
    {
      \begin{cvitems}
        \item {Analyse du besoin des chercheurs et développement en autonomie d’une application fonctionnelle pour la recherche d’articles scientifiques.}
        \item {Automatisation d’un workflow de collecte de métadonnées en Python, de l’exécution planifiée par tâches cron à la transmission des résultats par webhooks.}
        \item {Orchestration d’un agent IA multi-étapes sans framework dédié, avec l’API OpenAI, des prompts adaptés au cas d’usage et la gestion des appels d’outils, de l’état et du contexte.}
        \item {Intégration des composants via des API REST Python/FastAPI documentées avec OpenAPI, avec stockage dans MongoDB et interface React/Next.js pour les chercheurs.}
        \item {Mise en œuvre d’un pipeline RAG avec chunking, embeddings, indexation dans ChromaDB et recherche sémantique.}
        \item {Évaluation de plusieurs LLM selon la réussite de la recherche et la latence. Conteneurisation de composants avec Docker et mise en place d’un pipeline CI/CD sur GitHub.}
      \end{cvitems}
    }

  \cventry
    {Stagiaire ingénieur de recherche en apprentissage par renforcement}
    {Université Karunya}
    {Coimbatore, Inde}
    {Juin 2025 -- Septembre 2025}
    {
      \begin{cvitems}
        \item {Développement d’un modèle d’apprentissage par renforcement pour la navigation autonome d’un robot hospitalier, puis déploiement et tests sur Raspberry Pi.}
      \end{cvitems}
    }
\end{cventries}
